% Source: sources/returns/S6_04_MAPS_MANIFOLDS/sections/manifolds.tex:399–449
\paragraph{Theorem D.1.7 (Proved: dependence of the smooth type and the count).}
Within the multiplicity-$(3,4)$ construction, the oriented
diffeomorphism type of $X_p$ depends exactly on $d=|p|$.
For each $d$ coprime to six it carries a punctured-disc family of
pairwise non-biholomorphic complex structures, parametrized by
$W=e^{-\pi i(c+k_2)}$, $0<|W|<216$. At each fixed nonzero admissible
linearization scale $u=e^{-2\pi ic}$ it carries exactly two of these
biholomorphism classes. Its universal covers likewise have one
oriented diffeomorphism type for each $d$. This statement identifies
dependence on the construction data. The complete Zhubr classifying
invariant and a canonical cyclic action are supplied in D.1.9--D.1.12.
\begin{proof}
We use the all-twist comparison R1, whose sufficient step is the
integral norm equation and analytic torsor gluing in S6\_02,
B.2.1. Specifically, equality of $|p|$ and of $W$ gives
$p'=ap$, $a=\pm1$, and $c'-c\equiv k_2-k'_2\pmod2$.
The finite norm equations have integral solutions; the remaining
cusp defect is corrected in its entire rank-two quotient lattice,
and $H^1(\mathcal O\oplus\mathcal O(-1))=0$ glues the translations.
The cusp map preserves the prescribed fan. These are the steps that
give actual biholomorphisms, including every discrete twist, rather
than equality of abstract numerical data.

For fixed discrete data, varying $c$ along any compact path in the
admissible half-plane gives a proper smooth family. Here the
properness argument of S6\_02 A.3.1 uses a bounded correction
$C(t,c)$ in $Z=sB_0+C$: on a compact parameter path the effective
logarithmic lattice is uniformly close to $B_0$, so finitely many
closed toric polydiscs provide compact representatives. The finite
free quotients and the ordinary torus parallelepipeds are uniform
as well. All charts are products with the parameter before their
free identifications, so projection is a submersion. The same proof
does not use $p=\pm1$, and therefore applies to the present data.
Integrating a horizontal vector field over the compact path gives
an orientation-preserving diffeomorphism of its end fibres.

Given two data with the same $d$, first change one parameter by an
integer to align its value of $c+k_2$ with that of the other
discrete datum; use R1's biholomorphism at this aligned parameter,
and connect the resulting parameter to the desired one inside the
admissible half-plane. This proves independence of every datum
except $d$. Distinct $d$ give nonisomorphic fundamental groups.
Every resulting diffeomorphism lifts to an oriented diffeomorphism
of universal covers; different $d$ are distinguished upstairs by
Euler characteristic $2d$.

R1 makes the map to $W$ an exact biholomorphic classification.
Since $u=W^2$, its fibre at fixed $u\ne0$ is precisely the two
points $W$ and $-W$. No two of these are equal. The asserted count
and the punctured-disc family follow.
\end{proof}
